%-------------------------------%
%  Author: Alessandro Sciarra   %
%    Date: 08 Jun 2026          %
%-------------------------------%

\begin{exercise}[Inspirational]{Concurrency vs parallelism: The CPU illusion}
    In this exercise you will explore how Bash background execution interacts with the operating system scheduler and why ``starting many processes'' is not the same as ``running faster''.

    You are given the following function:
    \begin{lstlisting}[style=MyBash, numbers=none, emph={[2]{i, x, iterations}}]
        function Burn_Cpu()
        {
            local iterations=${1:-500000} x=0
            for ((i=0; i<iterations; i++)); do
                (( x += i ))
            done
        }
    \end{lstlisting}
    This function performs CPU work only (no I/O, no sleep).

    \begin{enumerate}
        \item Run the function once and measure its execution time using \bash|time Burn_Cpu|.
              Adjust the number of the function argument to make it run for a few seconds.
              Repeat the measurement a few times.
              Is the result stable?
        \item Determine the number of available CPU cores using \bash|nproc| (or equivalent on your OS).
        \item Define the following utility function:
              \begin{lstlisting}[style=MyBash, numbers=none, emph={[2]{i, n_concurrent}}]
                  function Time_Concurrent_Burn()
                  {
                      local n_concurrent=${1:-1}
                      echo "Timing ${n_concurrent} concurrent process(es):"
                      time {
                          for ((i=0; i < n_concurrent; i++)); do
                              Burn_Cpu &
                          done
                          wait
                      }
                  }
              \end{lstlisting}
              Run it with 2 processes and interpret the result.
        \item Run \bash|Time_Concurrent_Burn $(nproc)|.
              Compare the result with the single-process execution.

        \item Repeat the previous run, but now use twice the number of CPU cores:
              \begin{lstlisting}[style=MyBash]
                  Time_Concurrent_Burn $(( $(nproc)*2 ))
              \end{lstlisting}
              Observe what happens to \texttt{real}, \texttt{user}, and \texttt{sys} time.
        \item \textbf{Optional:}
              Write a \bash|Scan_Times| function to create a summary table of times.
              Remember that the Bash \bash|time| keyword prints its output to standard error.
              \begin{lstlisting}[style=MyBash]
                 |+ N           real           user            sys
                  1
                  2
                  3
                  4
                  5
                  6
                  7
                  8
                  9
                 10
                 12
                 16+|
              \end{lstlisting}
        \item Answer the following questions:
              \begin{itemize}[nosep]
                  \item Why can the total \texttt{user} time be larger than \texttt{real} time?
                  \item Why does adding more background processes not always reduce execution time?
                  \item Does \bash|&| guarantee parallel execution?
                        If not, what does it guarantee?
                  \item What role does the operating system scheduler play in these experiments?
              \end{itemize}
    \end{enumerate}
\end{exercise}